% Append-only section for the AION technical record.
% Required packages: amsmath, booktabs, array, url.

\section{Qwen3-30B Quality Recovery and Direct-I/O Cold-Start Programme}
\label{sec:qwen30b-quality-dio-2026-09-09}

\subsection{Scope and preserved baseline}

This section records the experiments completed after promotion of the
selective-Q2-expert Qwen3-30B-A3B execution path.  The preserved baseline is
the 10,516,205,056-byte (9.79~GiB) candidate with SHA-256
\texttt{7b1165eef9a868954a3fc9bdc31e3f454a04ebd86b14040945074c99edce5bac}.
It retains Q4\_K\_M policy for non-expert tensors, stores the MoE gate, up and
down expert banks as Q2\_K, loads all 48 layers on Metal, sustains more than
50 generated tokens/s through replicated 64-, 128- and 256-token tests, and
runs at approximately 61--63 tokens/s when fully warm.  These resident rates
must not be represented as continuous SD streaming rates: the SD card is the
durable cold warehouse and the compact model executes from unified memory
after loading.

Two independent tracks followed:
\begin{enumerate}
  \item a changed-weight, quality-gated mixed-precision track intended to
        recover broader model quality without losing residency; and
  \item a runtime-only direct-I/O track intended to reduce SD cold-start time
        without changing model arithmetic.
\end{enumerate}

\subsection{Heuristic edge-layer mixed precision}

The first exploratory candidate restored Q3\_K expert precision at layers
0--7 and 40--47, retained Q2\_K experts at layers 8--39, and retained the
Q4\_K\_M base policy for non-expert tensors.  This allocation was explicitly
declared as an architectural heuristic rather than a measured sensitivity
claim.  The predeclared plan SHA-256 was
\texttt{4b6912d12ef983e71dbe0db307095758541c7452522560fe3640e0a47eb12542}.

The temporary internal-disk candidate was 11,497,672,192 bytes
(10,959.34~MiB), with SHA-256
\texttt{1cec6825acb0a912a449ac0f5aacfc376ca4dce548d1317b741a3253438287f4}.
It was not copied to SD because the evidence volume had only approximately
9~GiB free.  It loaded with all layers on Metal and reached 58.13 tokens/s in
a narrow 16-token smoke test.

\begin{table}[ht]
\centering
\begin{tabular}{lccc}
\toprule
Evaluation & Q2 baseline & Edge candidate & Gate \\
\midrule
Structured strict & 6/8 & 7/8 & pass \\
Structured semantic & 7/8 & 8/8 & pass \\
Structured parseable & 8/8 & 8/8 & pass \\
Multiple choice & 11/12 & 10/12 & fail \\
\bottomrule
\end{tabular}
\caption{Frozen quality result for the heuristic edge-layer candidate.}
\end{table}

The structured cohort was repeated twice with identical raw and parsed
outputs.  The only strict miss used \texttt{gram} instead of the schema-exact
\texttt{g}, while returning the correct value.  However, the separate frozen
multiple-choice cohort regressed from 11/12 to 10/12.  The candidate was
therefore marked \textbf{NOT\_PROMOTED}; the predeclared replicated 256-token
stage was not run after the quality failure.  The canonical result SHA-256 is
\texttt{d47d6e0c65522c31f6c290c31b4cdba342c52d4c7a97a7860a7cf9d483eb0475}.

\subsection{Complete eight-layer sensitivity screen}

The heuristic result motivated a predeclared sensitivity screen.  The 48
transformer layers were divided into six independent eight-layer bands.  In
each candidate, exactly one band's expert gate, up and down tensors used
Q3\_K; all other expert tensors used Q2\_K.  A band was sensitivity-positive
only if it simultaneously achieved at least 7/8 strict and 7/8 semantic on
the structured cohort, 8/8 parseable and repeatable output, at least 11/12 on
the multiple-choice cohort, complete Metal residency, and at least 50
tokens/s on the smoke test.  The plan SHA-256 is
\texttt{90ddbde0c57add79285989d2d7b34562d5e8d4b23f719d9b00a9b9b08f3b5e7c}.

\begin{table}[ht]
\centering
\small
\begin{tabular}{lcccc}
\toprule
Band & Layers & Structured strict & MCQ & Result \\
\midrule
B0 & 0--7   & 6/8 & 10/12 & fail \\
B1 & 8--15  & 6/8 & 11/12 & fail \\
B2 & 16--23 & 7/8 & 10/12 & fail \\
B3 & 24--31 & not run & not run & Metal OOM \\
B4 & 32--39 & not run & not run & Metal OOM \\
B5 & 40--47 & 7/8 & 9/12  & fail \\
\bottomrule
\end{tabular}
\caption{Independent eight-layer Q3\_K expert-restoration screen.}
\end{table}

B0, B1, B2 and B5 produced repeatable, parseable structured outputs.  B0's
narrow smoke test measured 59.10 tokens/s.  B2 and B5 improved structured
results but regressed multiple-choice performance.  B1 preserved the Q2
multiple-choice score but did not improve the structured strict score.  B3 and
B4 encountered Metal command-buffer out-of-memory errors during initialisation.
B4 failed repeatedly even with one service slot, context 512, batch 128 and
microbatch 64.  Because those candidates were tested late in a
memory-intensive ordered run, temporal memory pressure remains a confound and
the failures cannot be attributed solely to layer position.

No band passed every frozen gate.  No band and no combination was promoted.
The screen demonstrates that precision response is non-monotonic and
task-dependent: restoring precision can repair one evaluation family while
damaging another.  Combining individually failing bands would be speculation,
not sensitivity-guided allocation.  The canonical result SHA-256 is
\texttt{d05d7dbf00c573c9164f3460477274b2a6b8c3204d97c4d5129db1bdfabf621a}.
The result was recorded in isolated commit \texttt{cb1b6fe9}.

\subsection{Direct-I/O narrow cold-start experiment}

The next track changed only the model-loading mechanism.  The question was
whether llama.cpp direct I/O could reduce cache-displaced SD startup without
changing output or resident performance.  The plan fixed a minimum 15\%
readiness reduction, identical deterministic token text and top-10
log-probability payload, at least 50 resident tokens/s, successful health, and
no Metal out-of-memory error.  Its SHA-256 is
\texttt{fc17589cbd0e47b1636a4cbeadee0b1db3fffd352261d49adabf56b30f8e0569}.

Immediately before the DIO start, the preserved 18,556,685,824-byte Q4 control
was read sequentially from the SD card to \texttt{/dev/null}.  The read took
194.096~s at 95.61~MB/s decimal.  This strongly displaced Q2 pages but is not
claimed as a kernel-certified complete cache flush.

\begin{table}[ht]
\centering
\begin{tabular}{lrr}
\toprule
Condition & Ready time (s) & Relative result \\
\midrule
Automatic mmap control & 213.088 & control \\
Direct I/O candidate & 116.219 & 45.459\% lower \\
\bottomrule
\end{tabular}
\caption{Initial cache-displaced service-readiness result.}
\end{table}

Direct I/O saved 96.869~s and made this measured start 1.8335 times as fast.
The fixed 32-token probe returned identical text, token sequence and top-10
log-probability payload.  Resident probe performance was 60.010 tokens/s for
DIO and 59.491 tokens/s for the preceding control.  All Metal layers loaded,
health became ready, and no out-of-memory error occurred.  The result passed
every narrow gate, while the document explicitly required balanced
replication before a general cold-start claim.  Its result SHA-256 is
\texttt{d226029a29310ca438f6796e9789f1a1599de81732841fa5f22a5289ce1bb543},
recorded in commit \texttt{510f23d8}.

\subsection{Balanced DIO replication}

A separate plan fixed the balanced order auto/DIO/DIO/auto.  Each arm first
read the full 17.28~GiB Q4 control from SD, then started the same Q2 model with
identical Metal, context, slot, thread, polling and attention settings.  The
plan SHA-256 is
\texttt{b81bb139e49a26dc0d500a41a582142a10e2ff1e4c0fc445ce503350068bd5ef}.

\begin{table}[ht]
\centering
\begin{tabular}{lrrr}
\toprule
Arm & Mode & Ready time (s) & First probe (tokens/s) \\
\midrule
1 & auto & 221.030 & 24.038 \\
2 & DIO  & 117.353 & 58.473 \\
3 & DIO  & 116.392 & 58.356 \\
4 & auto & 232.056 & 24.051 \\
\bottomrule
\end{tabular}
\caption{Balanced cache-displaced startup and immediate first-request rate.}
\end{table}

Median readiness was 226.543~s for automatic mmap and 116.873~s for DIO.
Direct I/O reduced median readiness by 48.410\%, a 1.9384-fold speedup.  Both
DIO arms were faster than both automatic arms.  All arms became healthy, none
reported Metal out-of-memory, and every completion produced identical text and
an identical top-10 log-probability payload.  The canonical log-probability
SHA-256 for every arm is
\texttt{1fe300c9fb23f5314920617a1a61980a8d27934aa4bc860752f3300ff445919b}.

The audit preserved one important predeclared-gate defect.  The plan required
\emph{every arm}, including controls, to exceed 50 tokens/s on the first
post-ready probe.  Both automatic arms delivered approximately 24 tokens/s,
while both DIO arms exceeded 58 tokens/s.  The literal all-arm gate therefore
failed and was not relaxed.  Formal promotion under this exact ABBA plan was
withheld, even though the DIO mechanism and direction replicated strongly.

The measurements support the hypothesis that mmap can announce readiness
while expensive first-use page work remains, whereas DIO supplies both earlier
readiness and immediately fast inference.  This is a measured systems
inference, not yet a causal decomposition.  A subsequent experiment must
separately measure process-start-to-first-completed-response and second-response
warm performance, with candidate acceptance separated from control
characterisation.

The balanced result SHA-256 is
\texttt{f9f064f18d6dd8ab327bfdf8aabf007c58996ea2e05dc2b9793572a374693498},
recorded in isolated commit \texttt{ac5939bf}.

\subsection{Current verified position}

The mixed-precision programme did not produce a quality-safe replacement and
was stopped without weakening its gates.  The original selective-Q2 model
therefore remains the resident quality-gated candidate.  The direct-I/O
programme, by contrast, produced a replicated and order-balanced systems
advance: approximately half of cache-displaced service-readiness time was
removed, and immediate first inference exceeded 58 tokens/s instead of the
automatic arms' approximately 24 tokens/s, with an identical deterministic
probability trace.

The precise claim is consequently:
\begin{quote}
A 30.532-billion-parameter sparse model stored on an ordinary SD card can be
loaded fully onto the existing 18~GB M3 Pro using AION's compact resident
warehouse.  Direct-I/O loading reduced measured cache-displaced readiness from
a 226.54-second balanced-control median to 116.87 seconds, while preserving the
tested token and top-log-probability trace.  After loading, generation occurs
from unified memory, not by streaming 60 tokens/s directly from SD.
\end{quote}

Generalisation remains open.  The next gates are larger independent startup
replication, process-start-to-first-response measurement, broader blinded model
quality evaluation, and testing on additional commodity Macs and storage
devices without altering the frozen evidence or claim boundary.

